% Chapter: one theorem behind every half (Euler's theorem for homogeneous functions)
\chapter{One theorem behind every half}\label{ch:euler_halves}

\section*{Overview}
A factor of one half appears at almost every place in this book: the kinetic share of a bound atom, star or cluster
(Chapter~\ref{ch:virial_law}), the share $T_HS$ of a black hole's mass (the Smarr relation), the coupling $\beta_m=\Omega_m/2$ at the cosmic horizon
(Chapter~\ref{ch:virial}), the energy a transistor gate stores out of what it draws (Chapter~\ref{ch:synthesis}) and the equal channel weights of
Koide's relation (Chapter~\ref{ch:koide}). This chapter shows that each of these exact halves is one theorem, Euler's theorem for homogeneous
functions, applied at one of two degrees, and states what that settles and what it leaves open. \derived

\section{The theorem}\label{sec:eh_theorem}
A function is homogeneous of degree $k$ when scaling every argument by $\lambda$ scales the function by $\lambda^k$. Euler's
theorem~\cite{Goldstein} then gives
\begin{equation}
f(\lambda\mathbf x)=\lambda^kf(\mathbf x)\quad\Longrightarrow\quad \mathbf x\cdot\nabla f=k\,f .\label{eq:eh_euler}
\end{equation}
\derived{} (standard). Every half below is this identity at a degree fixed by geometry, never by a fit.

\section{Degree $-1$: bound systems and Einstein's equations}\label{sec:eh_branch1}
A potential homogeneous of degree $k$ gives the virial theorem $2\langle K\rangle=k\langle V\rangle$ (Eq.~\eqref{eq:vl_euler}). Gravity and the
Coulomb force have $k=-1$ because Gauss's law in three spatial dimensions makes the field of a point source fall as $1/r^2$, so
$\langle K\rangle=\tfrac12|\langle V\rangle|$ for atoms (the quantum form by scaling, Fock~\cite{Fock1930}), stars and clusters of galaxies
(Table~\ref{tab:virial_domains}). \derived{} General relativity carries the same branch: for a stationary system the Tolman mass equals the ADM mass,
and that equality reduces to $2\langle K\rangle+\langle V\rangle=0$ in the weak field (Section~\ref{sec:vl_gr})~\cite{Tolman1930,GourgoulhonBonazzola1994}.
\derived{} The three-way cluster test reads this branch directly: in general relativity the hydrostatic and lensing masses of a relaxed cluster are
one mass, so a lensing-to-hydrostatic ratio that falls with redshift as $1/\mu(z)$ says that matter and light do not feel the same mass
(Chapter~\ref{ch:threeway}). \interp

\section{Degree 2: charge, horizons and modes}\label{sec:eh_branch2}
\paragraph{The capacitor.} The energy stored on a capacitor, $U=Q^2/2C$, is homogeneous of degree 2 in the charge. Charging it from a supply at the
final voltage $V=Q/C$ draws $QV=Q\,\partial U/\partial Q$, and Eq.~\eqref{eq:eh_euler} gives
\begin{equation}
Q\frac{\partial U}{\partial Q}=2U\quad\Longrightarrow\quad\frac{U}{QV}=\frac12 ,\label{eq:eh_cap}
\end{equation}
whatever the resistance. \derived{} This is the half that a transistor gate stores of what it draws (Chapter~\ref{ch:synthesis}).

\paragraph{The horizons.} The Bekenstein--Hawking entropy is proportional to the horizon area~\cite{Bekenstein1973,Hawking1975}, and for a
Schwarzschild black hole the area is proportional to $M^2$, so $S(M)$ is homogeneous of degree 2. With the first law $dMc^2=T_H\,dS$,
\begin{equation}
M\frac{dS}{dM}=2S\quad\Longrightarrow\quad T_HS=\frac{c^2}{dS/dM}\,S=\tfrac12Mc^2 ,\label{eq:eh_area}
\end{equation}
which is the Smarr relation and the Komar mass of Eq.~\eqref{eq:vl_komar} read as Euler's theorem. \derived{} At the cosmic horizon of radius $R$
the energy inside, $E=c^4R/2G$, grows as $R$ and the entropy as $R^2$, so $S(E)$ is again homogeneous of degree 2: the first-law temperature is
\begin{equation}
\frac{dE}{dS}=\frac{T_{GH}}{2},\qquad \frac{dE}{dS}\,S=\tfrac12E,\qquad T_{GH}S=E ,\label{eq:eh_dS}
\end{equation}
with $T_{GH}=\hbar c/2\pi k_BR$. \derived{} The half of the cosmic horizon is the same degree-2 identity as the black hole's.

\paragraph{Modes.} A mode whose energy is quadratic in its amplitude has $\langle K\rangle=\langle V\rangle$, and its orbit average keeps half its
peak, $\langle(y\cos\phi)^2\rangle=y^2/2$ (Eq.~\eqref{eq:ko:k2}). \derived{} On Koide's charge orbit this identity gives the weight of the
harmonic channel, $y^2/2$ (Eq.~\eqref{eq:ko:weights}). It does not give the equality of the two weights: Euler's theorem fixes degrees, not the
ratio of two channels of the same degree, and Koide's ratio is homogeneous of degree 0 in $(x,y)$, $Q=\tfrac13+y^2/6x^2$, so no normalization set
by the area law reaches it. \derived{} The equal split $x^2=y^2/2$, read in Section~\ref{sec:ko:virial} as equal power in the zero-sequence and
balanced parts of the square-root masses, remains the conjecture of Step~5. \conjecture

\section{What the shared root settles}\label{sec:eh_settles}
\begin{table}[htbp]\centering\small
\caption{Each exact half in the book as Euler's theorem, Eq.~\eqref{eq:eh_euler}. The degree is fixed by geometry: $-1$ by Gauss's law in three
dimensions, 2 by an energy quadratic in charge or amplitude, or by an entropy quadratic in mass (the area law).}\label{tab:eh_halves}
\begin{tabular}{@{}llcll@{}}\toprule
place & homogeneous function & degree & the half & chapter\\\midrule
atom, star, cluster & potential $V(\mathbf r)$ & $-1$ & $\langle K\rangle=\tfrac12|\langle V\rangle|$ & \ref{ch:virial_law}\\
relativistic bound system & Tolman $=$ ADM mass & $-1$ (weak field) & $2\langle K\rangle+\langle V\rangle=0$ & \ref{ch:virial_law}\\
black hole & $S(M)\propto M^2$ & 2 & $T_HS=\tfrac12Mc^2$ & \ref{ch:virial_law}\\
cosmic horizon & $S(E)\propto E^2$ & 2 & $(dE/dS)\,S=\tfrac12E$ & \ref{ch:euler_halves}\\
capacitor (transistor gate) & $U(Q)\propto Q^2$ & 2 & stored $=\tfrac12$ drawn & \ref{ch:synthesis}\\
Koide's harmonic channel & mode energy $\propto y^2$ & 2 & $\langle(y\cos\phi)^2\rangle=\tfrac12y^2$ & \ref{ch:koide}\\\bottomrule
\end{tabular}
\end{table}
Each entry of Table~\ref{tab:eh_halves} is Euler's theorem; none uses a fitted number. \derived{} Two consequences follow. First, the half of the
capacitor and the half of a black hole are the same theorem at the same degree: energy quadratic in charge and entropy quadratic in mass. The
conjecture of Chapter~\ref{ch:synthesis} that the transistor half and the virial half are one quantity is narrowed to this: they are one theorem
at two degrees, which is derived; that they are one quantity is still a conjecture. \conjecture{} Second, the gravitational branch is not separate
from the thermodynamic one. Jacobson derived Einstein's equations from $\delta Q=T\,\delta S$ on local horizons~\cite{Jacobson1995}, so the
equations that carry the degree $-1$ half come from the same local first law that IAM's Law prices; read with IAM's Law, the half in each row is the
share written on, or held by, the surface nearest the event. \interp

\section{What it does not settle}\label{sec:eh_open}
The root is shared; no number is transferred between rows, as Chapter~\ref{ch:synthesis} states. \interp{} The horizon's half is a
four-dimensional fact. In $D$ spacetime dimensions the mass of a static black hole grows as $r^{D-3}$ and its entropy as $r^{D-2}$, so
\begin{equation}
\frac{T_HS}{Mc^2}=\frac{D-3}{D-2},\label{eq:eh_smarrD}
\end{equation}
which is $\tfrac12$ only for $D=4$ ($\tfrac23$ for $D=5$). \derived{} The half kept by a mode on a circle, $\langle\cos^2\phi\rangle=\tfrac12$, does not
depend on $D$ (on a sphere it is $\tfrac13$). The horizon half and the orbit half therefore coincide in four dimensions, with no shared cause shown.
\derived{} Three things stay open. The cell's floor
$1/(1+e^M)$ is Boltzmann's two-state ratio, not one of these halves, and the cell's held share of its drive, $\varphi=0.1464$, is not a half (it is the slow end of the same write rule whose sudden end is the half; Chapter~\ref{ch:iama}); that its
zero-drive limit equals Koide's $Q=2/3$ has no shared reason found. \openprob{} Koide's offset $\delta$ and scale $x$ are not fixed by the
theorem (Section~\ref{sec:ko:offset}). \openprob{} The proportionality of entropy to area is taken from Bekenstein and Hawking, not derived here.
\openprob

\section{Tests}\label{sec:eh_tests}
The degree $-1$ branch is tested in every bound system of Table~\ref{tab:virial_domains}, and across the sector split by the sign of the redshift
slope of the lensing-to-hydrostatic ratio (Chapter~\ref{ch:threeway}). \prediction{} The degree 2 reading of Koide's relation keeps its test:
$Q=2/3$ exactly fixes the $\tau$ mass (Section~\ref{sec:ko:tests}). \prediction

\section{Status}\label{sec:eh_status}
\begin{center}\small\begin{tabular}{p{0.66\textwidth}l}\toprule
statement & label\\\midrule
Euler's theorem, Eq.~\eqref{eq:eh_euler} & \derived\\
the capacitor half, Eq.~\eqref{eq:eh_cap}; the horizon halves, Eqs.~\eqref{eq:eh_area}--\eqref{eq:eh_dS} & \derived\\
the degree $-1$ half from Einstein's equations (Tolman $=$ ADM, weak field) & \derived\\
the transistor half and the virial half as one theorem at two degrees & \derived\\
the transistor half and the virial half as one quantity & \conjecture\\
Koide's harmonic weight $y^2/2$ as the degree 2 half on the charge orbit & \derived\\
Koide's equal split $x^2=y^2/2$ & \conjecture\\
the horizon half four-dimensional, Eq.~\eqref{eq:eh_smarrD} & \derived\\
the half as the share written on the nearest surface & \interp\\
the cell's floor and held share on this thread; Koide's $\delta$ and $x$ & \openprob\\\bottomrule
\end{tabular}\end{center}
